% Full title: Outgoing Cauchy Relations and Bounded-Domain Reduction
\section{Outgoing Cauchy Relations and Bounded-Domain Reduction}
\label{sec:cauchy-reduction}

\subsection{Weak traces and outgoing half-guide spaces}

Let $D$ be one of $\mathcal C^0,W^-,W^+$.  For a field
$u\in H^1_\beta(D)$ satisfying $\Delta u\in L^2(D)$ piecewise, the Dirichlet
trace on a port belongs to $H^{1/2}_\beta(\Gamma)$.  Its weak outward Neumann
trace $\gamma_{N,D}u\in H^{-1/2}_\beta(\Gamma)$ is defined by
\begin{equation}
  \pair{\gamma_{N,D}u}{\varphi}_{\Gamma}
  :=\int_D\left(\nabla u\cdot\nabla\overline V-k^2\rho\,u\overline V\right)
  \dd x\dd y,
  \label{eq:weak-Neumann}
\end{equation}
where $V\in H^1_\beta(D)$ has port trace $\varphi$, vanishes on any other
finite port, and is otherwise an arbitrary bounded lifting.  The weak PDE
makes \eqref{eq:weak-Neumann} independent of the lifting.  Horizontal terms
cancel by quasiperiodicity.

The unit normals pointing \emph{out of the center cell} are
\[
  \nu^{0,-}=-e_x\quad\text{on }\Gamma^- ,\qquad
  \nu^{0,+}=e_x\quad\text{on }\Gamma^+.
\]
They are opposite to the outward normals of the adjacent half-guides.  We use
this center-oriented convention everywhere.  For $u^0$ in $\mathcal C^0$ and
$w^\pm$ in $W^\pm$, define
\begin{align}
  \Tr^{0,\pm} u^0
  &:=\bigl(\gamma_Du^0,\gamma_{N,\mathcal C^0}u^0\bigr)
  \in\mathcal H_\Gamma^\pm,\label{eq:center-trace}\\
  \Tr^{0,\pm} w^\pm
  &:=\bigl(\gamma_Dw^\pm,-\gamma_{N,W^\pm}w^\pm\bigr)
  \in\mathcal H_\Gamma^\pm.
  \label{eq:halfguide-trace}
\end{align}
Consequently, equality of the two pairs means continuity of the field and of
the derivative in the same physical direction $\nu^{0,\pm}$.

Let $H^1_{\beta,0}(W^\pm;\Gamma^\pm)$ be the subspace with zero Dirichlet
trace at the finite port.  The square-integrable half-guide solution space is
\begin{equation}
  \mathcal U^\pm_{\mathrm{out}}(k,\beta):=
  \left\{w\in H^1_\beta(W^\pm):
  \int_{W^\pm}\nabla w\cdot\nabla\overline v
  =k^2\int_{W^\pm}\rho\,w\overline v
  \ \forall v\in H^1_{\beta,0}(W^\pm;\Gamma^\pm)\right\}.
  \label{eq:halfguide-outgoing}
\end{equation}
When $k^2\in I_\beta$, this $H^1$ condition selects the stable, decaying
solutions.  No pointwise radiation condition is added.

\begin{definition}[Outgoing Cauchy relations]
\label{def:outgoing-relations}
The outgoing Cauchy relation of the $\pm$ half-guide is
\begin{equation}
  \mathcal C^\pm_{\mathrm{out}}(k,\beta)
  :=\left\{\Tr^{0,\pm} w:w\in
  \mathcal U^\pm_{\mathrm{out}}(k,\beta)\right\}
  \subset\mathcal H_\Gamma^\pm.
  \label{eq:outgoing-relation}
\end{equation}
It is a linear subspace of complete Dirichlet--Neumann pairs; it is not assumed
to be the graph of an operator.  The Dirichlet and Neumann projections are
\begin{equation}
  \pi_D^\pm(f,g):=f,\qquad \pi_N^\pm(f,g):=g,
  \qquad (f,g)\in\mathcal H_\Gamma^\pm.
  \label{eq:trace-projections}
\end{equation}
\end{definition}

\begin{lemma}[Weak gluing across an artificial port]
\label{lem:weak-gluing}
Let $u_1$ and $u_2$ be $H^1_\beta$ weak Helmholtz solutions in two adjacent
Lipschitz subdomains.  If their Dirichlet traces agree and their outward weak
Neumann traces sum to zero on the common port, then their piecewise union is an
$H^1_\beta$ weak solution on the union.  Conversely, the restrictions of a
global weak solution satisfy these two matching conditions.
\end{lemma}
\status{background result.}
\begin{proofroadmap}
The Sobolev gluing theorem gives a global $H^1$ function from equality of the
Dirichlet traces.  Apply Green's identity in the two subdomains.  The only new
distributional term is the sum of the two outward conormal traces; it vanishes
by hypothesis.  This is precisely why the sign in \eqref{eq:halfguide-trace} is
required.  The argument uses full Cauchy data and never infers uniqueness from
zero Dirichlet trace alone.

See \path{research/planning/muller-cauchy/notes/trace-glue.md} and
the obligation named ``restriction--gluing isomorphism'' in
\path{research/planning/muller-cauchy/p-proofs.md}.  The weakest
acceptable conclusion is distributional gluing; stronger local regularity may
then be recovered away from corners and material interfaces.
\end{proofroadmap}
\begin{proof}
Let $D_1$ and $D_2$ be the two subdomains and let
$\Gamma:=\partial D_1\cap\partial D_2$ denote their common port.  We prove
the two directions separately.

\emph{Step 1: Sobolev gluing.}
Set $u=u_j$ almost everywhere in $D_j$.  Equality of the Dirichlet traces
implies $u\in H^1_\beta(D_1\cup D_2)$.  For completeness, this follows
locally by flattening $\Gamma$ and computing the distributional derivatives:
the derivative of the piecewise function is the piecewise derivative plus an
interface distribution whose density is the jump
$\gamma_Du_1-\gamma_Du_2$.  The latter is zero in
$H^{1/2}_\beta(\Gamma)$, so the distributional derivatives belong to $L^2$.
The quasiperiodic trace condition on the horizontal sides is inherited from
the two pieces.

\emph{Step 2: cancellation of the conormal traces.}
Take $v\in H^1_\beta(D_1\cup D_2)$, initially with bounded support if the
union is unbounded.  The weak Green identities in the two subdomains give
\begin{align*}
  &\sum_{j=1}^2\int_{D_j}
  \left(\nabla u_j\cdot\nabla\overline v
  -k^2\rho\,u_j\overline v\right)\dd x\dd y\\
  &\qquad=
  \pair{\gamma_{N,D_1}u_1+\gamma_{N,D_2}u_2}
  {\gamma_Dv}_{\Gamma}.
\end{align*}
All other boundary terms either vanish by the imposed test space or cancel
by quasiperiodicity.  The right-hand side is zero by hypothesis.  Density of
the bounded-support tests in the relevant global test space then shows that
$u$ satisfies the weak Helmholtz equation on the union.

\emph{Step 3: the converse.}
Let $u\in H^1_\beta(D_1\cup D_2)$ be a global weak solution and set
$u_j=u|_{D_j}$.  A global $H^1$ function has a single Dirichlet trace on the
interior Lipschitz interface, hence
$\gamma_Du_1=\gamma_Du_2$.  Applying the two weak Green identities and
subtracting the global weak equation yields
\[
  \pair{\gamma_{N,D_1}u_1+\gamma_{N,D_2}u_2}{\varphi}_{\Gamma}=0
  \qquad\forall\varphi\in H^{1/2}_\beta(\Gamma),
\]
where an arbitrary $\varphi$ is realized by a bounded lifting supported near
$\Gamma$.  Thus the two outward weak Neumann traces sum to zero in
$H^{-1/2}_\beta(\Gamma)$, as claimed.
\end{proof}

Define the bounded center solution space
\begin{equation}
\begin{split}
  \mathcal B^0(k,\beta):=\bigl\{u^0\in H^1_\beta(\mathcal C^0):\ &
  u^0\text{ satisfies the center weak Helmholtz equation},\\[-0.2em]
  &\Tr^{0,-}u^0\in\mathcal C^-_{\mathrm{out}}(k,\beta),\quad
  \Tr^{0,+}u^0\in\mathcal C^+_{\mathrm{out}}(k,\beta)\bigr\},
\end{split}
\label{eq:bounded-center-space}
\end{equation}
where the center weak equation is tested against functions in
$H^1_\beta(\mathcal C^0)$ with zero Dirichlet trace on both ports.  The
coefficient $\rho$ encodes the physical transmission conditions at $\Upsilon^0$.

\begin{theorem}[Bounded trace-relation reduction]
\label{thm:bounded-reduction}
Assume $k^2\in I_\beta$ and that the half-guide relations in
\eqref{eq:outgoing-relation} are closed.  Restriction to the center cell defines
a linear isomorphism
\begin{equation}
  \operatorname{Res}^0:\mathcal G(k,\beta)
  \longrightarrow\mathcal B^0(k,\beta),
  \qquad u\longmapsto u|_{\mathcal C^0}.
  \label{eq:restriction-isomorphism}
\end{equation}
The inverse is obtained by choosing half-guide fields with the prescribed
Cauchy pairs and gluing them to the center field.  In particular,
\begin{equation}
  \dim\mathcal G(k,\beta)=\dim\mathcal B^0(k,\beta).
  \label{eq:bounded-multiplicity}
\end{equation}
\end{theorem}
\status{adaptation of a known result.}
\begin{proofroadmap}
For the forward map, restrict a global guided mode to the three subdomains;
the half-guide restrictions lie in \eqref{eq:halfguide-outgoing} and yield the two
relation conditions.  If the center restriction vanishes, its complete Cauchy
data vanish, and weak unique continuation gives a zero global field.

Conversely, choose half-guide preimages of the two relation pairs and apply
\Cref{lem:weak-gluing}.  Independence of the choices follows from uniqueness
for a half-guide field with zero complete Cauchy data.  Exponential decay is
then inherited from the half-guide spaces in the common gap.  For identical
half-guides this reduces to the DtN reduction of
\cite[Theorem~4.5]{Fliss2013}; the present asymmetric statement instead uses
the two Cauchy relations and the gluing argument independently at each end.

The internal audit is
\path{research/planning/muller-cauchy/r-draft-thm.md};
the detailed proof obligation is in
\path{research/planning/muller-cauchy/p-proofs.md}.  If closedness
or half-guide uniqueness fails, the weakest acceptable replacement is an
isomorphism after quotienting the half-guide zero-Cauchy solution space.
\end{proofroadmap}
\begin{proof}
We first record the uniqueness fact used below.  Under the geometry of
\Cref{sec:functional-setting}, a global piecewise Helmholtz solution which
vanishes on a nonempty open subset of the matrix phase vanishes everywhere.
Indeed, the matrix region obtained by removing the mutually separated compact
inclusions from the quasiperiodic strip is connected.  In that region the
coefficient is constant, so interior elliptic regularity and the Helmholtz
equation make the solution real analytic; vanishing on one open subset
therefore implies vanishing throughout the matrix region.  On every inclusion
boundary both the Dirichlet trace and the normal derivative consequently
vanish.  Extending the interior field by zero across a boundary patch and
using \Cref{lem:weak-gluing} gives a constant-coefficient Helmholtz solution
which vanishes on one side of that patch.  Analytic continuation makes it zero
inside the inclusion as well.  Repeating this argument proves the asserted
uniqueness on the whole strip.  The same argument applies to a half-guide
solution with zero complete Cauchy data after it is extended by zero across
its finite port.

\emph{Step 1: the restriction map is well defined and injective.}
Let $u\in\Gspace$.  Testing the global weak equation with zero extensions of
functions in $H^1_{\beta,0}(W^\pm;\Gamma^\pm)$ shows that
$u|_{W^\pm}\in\mathcal U^\pm_{\mathrm{out}}(k,\beta)$.  The center
restriction satisfies the center weak equation.  The converse part of
\Cref{lem:weak-gluing}, together with the sign convention in
\eqref{eq:halfguide-trace}, yields
\[
  \Tr^{0,\pm}(u|_{\mathcal C^0})
  =\Tr^{0,\pm}(u|_{W^\pm})
  \in\mathcal C^\pm_{\mathrm{out}}(k,\beta).
\]
Hence $\operatorname{Res}^0u\in\mathcal B^0(k,\beta)$.  If
$\operatorname{Res}^0u=0$, then $u$ vanishes on the open center matrix
region, and the uniqueness argument above gives $u=0$.  Thus
$\operatorname{Res}^0$ is injective.

\emph{Step 2: surjectivity by weak gluing.}
Take $u^0\in\mathcal B^0(k,\beta)$.  By the definition of the two outgoing
relations there are fields $w^\pm\in\mathcal U^\pm_{\mathrm{out}}(k,\beta)$
such that
\[
  \Tr^{0,\pm}w^\pm=\Tr^{0,\pm}u^0.
\]
The first components give equality of the Dirichlet traces.  By
\eqref{eq:center-trace}--\eqref{eq:halfguide-trace}, equality of the second
components is exactly
\[
  \gamma_{N,\mathcal C^0}u^0+\gamma_{N,W^\pm}w^\pm=0
  \quad\text{on }\Gamma^\pm.
\]
Applying \Cref{lem:weak-gluing} at both ports therefore produces a field
\[
  u=\begin{cases}
    w^-&\text{in }W^-,\\
    u^0&\text{in }\mathcal C^0,\\
    w^+&\text{in }W^+
  \end{cases}
  \in H^1_\beta(B)
\]
which satisfies the global weak eigenvalue equation.  The closed-form
characterization of $A(\beta)$ recorded in \Cref{prop:weak-strong} gives
$u\in\Gspace$, and by construction $\operatorname{Res}^0u=u^0$.

\emph{Step 3: independence of the half-guide preimages.}
If $\widetilde w^\pm$ is another choice with the same complete Cauchy pair,
then $z^\pm=w^\pm-\widetilde w^\pm$ has zero Dirichlet and center-oriented
Neumann traces at $\Gamma^\pm$.  Extend $z^\pm$ by zero through the center.
The gluing lemma makes this extension a global weak solution which vanishes
on a nonempty open set; the uniqueness argument at the start of the proof
implies $z^\pm=0$.  Hence the glued field is independent of all choices and
defines a linear inverse of $\operatorname{Res}^0$.

We have proved the isomorphism \eqref{eq:restriction-isomorphism}; equality of
dimensions in \eqref{eq:bounded-multiplicity} follows immediately.  The
assumed closedness of the two relations ensures that the two trace constraints
defining $\mathcal B^0(k,\beta)$ are closed; the algebraic restriction--gluing
argument itself uses only their defining trace-range property.
\end{proof}

\subsection{DtN charts and generalized Bloch traces}

Whenever $\pi_D^\pm|_{\Coutpm}$ is bijective onto
$H^{1/2}_\beta(\Gamma^\pm)$, define the DtN operator
\begin{equation}
  \Lambda^\pm(k,\beta):=
  \pi_N^\pm\bigl(\pi_D^\pm|_{\Coutpm}\bigr)^{-1}.
  \label{eq:DtN-definition}
\end{equation}
A failure of this inverse is a failure of the Dirichlet chart, not necessarily
of the Cauchy relation.  A Robin projection gives an alternative RtR chart as
in \cite{FlissKlindworthSchmidt2015}.

We next define the translation objects used below.  Extend the coefficient of
each half-guide periodically to a perfect two-sided strip.  Let
$\mathscr T^\pm(k,\beta)$ be the Hilbert space of center-oriented Cauchy traces
on $\Gamma_0^\pm$ of $H^1_{\mathrm{loc},\beta}$ solutions in that perfect
strip, equipped with a fixed local solution norm.  Unique continuation makes
the one-cell trace translation well defined on this solution-trace space:
\begin{equation}
  T^\pm(k,\beta)e:=\Tr^{0,\pm} v\big|_{\Gamma_1^\pm},
  \label{eq:translation-operator}
\end{equation}
where $v$ is the half-guide solution having trace $e$ on $\Gamma_0^\pm$, and the
trace at $\Gamma_1^\pm$ is transported back to the reference port by periodic
identification.  Boundedness of $T^\pm$ and equivalence of the local trace norm
with the solution norm are part of the translation-framework hypotheses.

\begin{lemma}[A common essential gap excludes propagating multipliers]
\label{lem:gap-unit-circle}
Assume the perfect-guide Floquet transform in the transverse periodic direction
is valid and that nonzero solution traces correspond to nonzero Bloch
solutions.  If $k^2\in I_\beta$, then
\begin{equation}
  \sigma(T^\pm(k,\beta))\cap\Sone=\varnothing.
  \label{eq:no-unit-multiplier}
\end{equation}
\end{lemma}
\status{background Floquet consequence.}
\begin{proof}
If a unit-modulus multiplier existed, write it as
$\lambda=e^{\ii\theta L^\pm}$ with real transverse quasimomentum $\theta$.
The corresponding nonzero solution would satisfy the one-cell transverse
Bloch condition and would place $k^2$ in the $\theta$-fiber spectrum of the
perfect guide.  The Floquet direct-integral decomposition would then give
$k^2\in\sigma(A^\pm(\beta))\subset\Sigma_{\mathrm{ess}}(\beta)$, contrary to
\eqref{eq:strict-gap}.  This is the spectral meaning of the absence of
propagating half-guide channels in a common essential gap.
\end{proof}

A generalized Bloch trace chain for a multiplier $\lambda\in\Complex$ is a finite
sequence $e_0,\ldots,e_{m-1}\in\mathscr T^\pm$ satisfying
\begin{equation}
  (T^\pm-\lambda I)e_0=0,
  \qquad (T^\pm-\lambda I)e_\ell=e_{\ell-1},
  \quad 1\leq\ell<m.
  \label{eq:Jordan-chain}
\end{equation}
The corresponding fields are generalized Bloch modes.  If
$\sigma(T^\pm)\cap\Sone=\varnothing$, choose a contour $\Gamma_s$ enclosing
precisely the multipliers with $|\lambda|<1$ and define
\begin{equation}
  P_s^\pm:=\frac{1}{2\pi\ii}\int_{\Gamma_s}(zI-T^\pm)^{-1}\dd z,
  \qquad \mathscr E_s^\pm:=\ran P_s^\pm.
  \label{eq:stable-projection}
\end{equation}
The subspace $\mathscr E_s^\pm$ is the stable spectral subspace.

Under a Riesz-basis hypothesis, collect all stable chains in a coefficient
Hilbert space
\begin{equation}
  \mathcal K_s^\pm:=
  \ell^2\bigl(\{(j,\ell): |\lambda_j^\pm|<1,
  0\leq\ell<m_j^\pm\}\bigr).
  \label{eq:stable-coefficients}
\end{equation}
Let
\begin{equation}
  \mathcal E^\pm(k,\beta):\mathcal K_s^\pm
  \longrightarrow\mathcal U^\pm_{\mathrm{out}}(k,\beta)
  \label{eq:field-synthesis}
\end{equation}
be the Riesz synthesis of the associated generalized fields, and define the
complete Cauchy-trace synthesis
\begin{equation}
  Q^\pm(k,\beta):=\Tr^{0,\pm}\mathcal E^\pm(k,\beta):
  \mathcal K_s^\pm\longrightarrow\mathcal H_\Gamma^\pm.
  \label{eq:trace-synthesis}
\end{equation}

\begin{proposition}[Cauchy relation, DtN graph, and stable traces]
\label{prop:relation-Bloch}
Assume the periodic half-guide and trace problem satisfy the hypotheses of the
generalized Floquet Riesz-basis theorem in
\cite[Theorem~2.2]{HohageSoussi2013}, and assume
$\sigma(T^\pm)\cap\Sone=\varnothing$.  Then
\begin{equation}
  \mathcal C^\pm_{\mathrm{out}}(k,\beta)=\ran Q^\pm(k,\beta),
  \label{eq:relation-stable-range}
\end{equation}
where the range is closed and the synthesis includes every stable Jordan
chain.  If in addition $\pi_D^\pm|_{\Coutpm}$ is bijective, then
\begin{equation}
  \mathcal C^\pm_{\mathrm{out}}(k,\beta)
  =\graph\Lambda^\pm(k,\beta).
  \label{eq:relation-DtN-graph}
\end{equation}
Ordinary Bloch eigenvectors may replace the chains only when the stable part
of $T^\pm$ is semisimple.
\end{proposition}
\status{known modal result with an adaptation required for the present transmission traces.}
\begin{proofroadmap}
First identify a unit-modulus multiplier with a real transverse Bloch
quasimomentum.  The common essential gap excludes such a perfect-guide Bloch
solution and provides separated stable and unstable Riesz contours.  Apply
\cite[Theorem~2.2 and Propositions~6.1--6.2]{HohageSoussi2013} to obtain a
Riesz basis of generalized modes and, when the trace problem is well posed, a
Riesz basis of their traces.  The spectral-decomposition route in
\cite[Sections~3--5]{Zhang2021} and the modal DtN construction in
\cite{Zhang2023} provide complementary formulations.

The nontrivial adaptation is to verify the coefficient, interface, and trace
isomorphism hypotheses for the piecewise-constant transmission half-guide used
here, including the two-component center-oriented Cauchy trace.  The graph
identity then follows immediately from \eqref{eq:DtN-definition}.  Consult
\path{research/planning/muller-cauchy/notes/bloch-trans.md}
and \path{research/planning/muller-cauchy/notes/outgoing.md}.
If a trace Riesz basis cannot be obtained, the weakest acceptable result is
the closed-range identity \eqref{eq:relation-stable-range} with coefficients
factored by $\ker Q^\pm$.
\end{proofroadmap}
\begin{proof}
It suffices to argue for one sign, denoted by
$\star\in\{-,+\}$, because the outward coordinate
$t^\star=s^\star(x-X^\star)$ puts both half-guides into the same right-going
half-strip form.

We first match the hypotheses of
\cite[Theorem~2.2 and Proposition~3.5]{HohageSoussi2013} in the order in
which they enter the argument.
\begin{enumerate}[label=\textup{(\arabic*)},leftmargin=*]
  \item Their scalar half-strip equation requires a positive, bounded,
  periodic coefficient.  After rescaling $t^\star$ by $L^\star$, this is
  precisely the periodic extension of \eqref{eq:material}; piecewise-constant
  transmission coefficients are allowed because the cited theorem assumes
  only an $L^\infty$ coefficient bounded away from zero.
  \item Their lateral boundary operator may be quasiperiodic.  The present
  $y$-quasiperiodicity is that case.  Any additional symmetry required by the
  cited theorem at its exceptional quasiperiodic parameters is included in
  the present proposition's explicit assumption that all of that theorem's
  hypotheses hold.
  \item Their outgoing solution space must be invariant under one-cell
  translation.  Here this is part of the translation-framework hypotheses.
  Moreover, $\sigma(T^\star)\cap\Sone=\varnothing$ removes all real-
  quasimomentum Floquet modes, so their radiating space
  $H^{1,+}$ reduces to the square-integrable $H^1$ solution space
  $\mathcal U^\star_{\mathrm{out}}(k,\beta)$.
  \item For the trace conclusion in part 2 of their theorem, the boundary
  problem for a trace
  $\tau=\theta_D\gamma_D+\theta_N\partial_{t^\star}$ must be well posed.
  This is exactly the ``trace problem'' part of the hypotheses assumed in the
  present proposition.  With the center-oriented convention,
  $\partial_{t^\star}$ is the second component of
  \eqref{eq:halfguide-trace}, so $\tau$ is a bounded projection of the present
  two-component Cauchy trace.
\end{enumerate}
Thus the solution-space and trace-space conclusions of the cited theorem,
and the Jordan-chain identification of its Proposition~3.5, apply without a
change of multiplier or orientation.

\emph{Step 1: identify the stable synthesis with all outgoing fields.}
The cited Riesz-basis theorem gives constants $a,b>0$ and a synthesis
isomorphism
\[
  \mathcal E^\star:\mathcal K_s^\star
  \longrightarrow\mathcal U^\star_{\mathrm{out}}(k,\beta)
\]
formed from the canonical generalized Floquet chains.  The unit-circle
exclusion implies that every mode in this $H^1$ solution space has a stable
multiplier and that every stable chain occurs in the synthesis.  Consequently,
using \eqref{eq:trace-synthesis} and then
\eqref{eq:outgoing-relation},
\[
  \ran Q^\star
  =\Tr^{0,\star}\ran\mathcal E^\star
  =\Tr^{0,\star}\mathcal U^\star_{\mathrm{out}}(k,\beta)
  =\mathcal C^\star_{\mathrm{out}}(k,\beta).
\]

\emph{Step 2: prove that the range is closed.}
Let $\mathcal Y_\tau^\star$ be the boundary space in which the well-posed
trace problem takes its data, and let
\[
  L_\tau^\star:\mathcal H_\Gamma^\star\to
  \mathcal Y_\tau^\star,
  \qquad L_\tau^\star(f,g)=\theta_Df+\theta_Ng.
\]
Part 2 of \cite[Theorem~2.2]{HohageSoussi2013} says that
$L_\tau^\star Q^\star$ is a Riesz synthesis isomorphism.  Hence, for some
$c_\tau>0$ and every $c\in\mathcal K_s^\star$,
\[
  c_\tau\lVert c\rVert_{\mathcal K_s^\star}
  \leq\lVert L_\tau^\star Q^\star c\rVert_{\mathcal Y_\tau^\star}
  \leq\lVert L_\tau^\star\rVert
  \lVert Q^\star c\rVert_{\mathcal H_\Gamma^\star}.
\]
The trace theorem and the upper Riesz estimate for $\mathcal E^\star$ also
give
\[
  \lVert Q^\star c\rVert_{\mathcal H_\Gamma^\star}
  \leq C\lVert c\rVert_{\mathcal K_s^\star}.
\]
Thus $Q^\star$ is bounded below as well as bounded.  In particular it is
injective and has closed range.  This proves all assertions in
\eqref{eq:relation-stable-range}.

\emph{Step 3: pass to the Dirichlet chart.}
Assume now that
$\pi_D^\star|_{\mathcal C^\star_{\mathrm{out}}}$ is bijective.  For every
$(f,g)\in\mathcal C^\star_{\mathrm{out}}$, definition
\eqref{eq:DtN-definition} gives
\[
  g=\pi_N^\star(f,g)
   =\pi_N^\star
    \bigl(\pi_D^\star|_{\mathcal C^\star_{\mathrm{out}}}\bigr)^{-1}f
   =\Lambda^\star f.
\]
Conversely, the inverse Dirichlet projection puts
$(f,\Lambda^\star f)$ in the relation for every
$f\in H^{1/2}_\beta(\Gamma^\star)$.  Hence
$\mathcal C^\star_{\mathrm{out}}=\graph\Lambda^\star$, which is
\eqref{eq:relation-DtN-graph}.

Finally, Proposition~3.5 of the cited paper identifies each root subspace with
finite Jordan chains for one-cell translation.  If a stable Jordan block has
size greater than one, its associated vectors are not eigenvectors and cannot
belong to the span of ordinary eigenvectors of that block.  Ordinary Bloch
eigenvectors can therefore replace all generalized chain vectors exactly when
the stable restriction of $T^\star$ is semisimple.
\end{proof}
